% FORGE algorithm figure. Vector, no external assets, no numbers: this states the method, so nothing
% in it is a result and nothing needs a hash pin.
%
% Composition: everything is centred on one vertical axis at x=7 and laid out in four even bands,
% so the figure reads as a stack rather than a scatter.
%
% Deliberately plain: no gradient ground, no glow, no dot lattice, no drop shadows, minimal corner
% rounding. Those effects are the tells of a generated-looking figure and none of them carried
% information here. Flat fills, hairlines and typography do the work instead.
%
% The four precursor-role colours are untouched: they are the key shared with the atlas and the
% generated-sample figure, so they are not a styling choice. What is compensated below is how
% heavily each one renders, not which colour it is.
%
% Claim discipline: exact replay and route closure are computational checks. The figure asserts
% neither experimental synthesis success nor biological activity.
\definecolor{roleamine}{HTML}{8A5AC2}
\definecolor{roleald}{HTML}{168C8C}
\definecolor{roleiso}{HTML}{E39532}
\definecolor{navybar}{HTML}{4C6780}
\definecolor{ground}{HTML}{F4F1E9}
\definecolor{groundcore}{HTML}{E7DCC5}
\definecolor{peri}{HTML}{55708A}
\definecolor{cyanacc}{HTML}{9C8062}
\definecolor{ink}{HTML}{22262F}
\definecolor{inkmuted}{HTML}{5C6474}

\newcommand{\plateref}{\tiny\color{white!84!navybar}}

\begin{tikzpicture}[
  font=\scriptsize,
  plate/.style={rounded corners=2pt, fill=navybar, text=white, inner xsep=10pt, inner ysep=6pt,
    minimum width=10.60cm, align=center, font=\footnotesize},
  statelab/.style={text=peri!85!black, font=\small},
  lead/.style={-{Stealth[length=3pt,width=3pt]}, draw=peri!75, line width=0.6pt},
]

\def\cx{7.00}      % the single vertical axis everything is centred on
\def\bandl{2.05}   % band block, left and right. The endpoints sit outside it, so these
\def\bandr{11.95}  % two values are what set the clearance around P and T.
\def\labx{2.20}    % label column. Set by the tightest pair in the block: "aldehyde" against
                   % the leftmost X_1 atom halo, and "amine" against its own, both ~0.32cm.
\def\inx{0.90}     % P, far enough left that "program in" clears the panel edge
\def\outx{13.05}   % T, mirrored. Not \px/\tx: those are loop variables further down.
\def\ya{0.05}      % trajectory band
\def\yaxis{2.60}   % sampling-time axis, above the molecules
\def\ypanel{-2.80} % coordinate panel
\def\bandh{0.35}   % band half-height. Holds the two-line stratum label; an atom halo is 0.17 of
                   % that, so both clear the band edge. At the original 0.155 the halos
                   % overflowed their own stratum and bled into the neighbouring ones.

%% ================================================================ ground
% A vignette rather than a panel: warm at the centre, fading to the page at the edges, so the
% figure sits on a field with no border to read as an arbitrary rectangle. The large corner
% radius plus the fade gives a lozenge rather than a box.
% Deliberately not an hourglass. In the transition-path figures this shape comes from, the waist
% is the transition state between two basins and the shape is the claim. Here the vertical axis
% is the role coordinate, so a waist would assert that roles merge as sampling proceeds, which
% is false: \Pi_P pins them at every step. The shape stays neutral and the strata carry meaning.
\shade[rounded corners=1.9cm, inner color=groundcore, outer color=white]
  (-0.15,-4.50) rectangle (14.10,4.46);
% Soft basins under the two endpoints, marking where the program enters and the route leaves.
\shade[inner color=peri!22, outer color=white] (\inx,\ya) circle (1.05);
\shade[inner color=cyanacc!22, outer color=white] (\outx,\ya) circle (1.05);


%% ================================================================ role strata
% The vertical direction is the role coordinate: one stratum per precursor role, and every atom is
% drawn inside its own stratum. This is what makes the background a coordinate system rather than a
% backdrop, and it is only legible because the sites are laid out by role instead of by chemical
% geometry.
% Three strata, not four. This axis is the role coordinate r_i, and r_i ranges over exactly the
% three Ugi-3CR precursor roles: ROLE_NAMES in forge/potency/annotations.py has three entries, and
% forge/model/ugi_joint_sparse_flow.py sizes role_embedding to len(ROLE_NAMES) -- the +1 slot is
% NULL_ROLE_INDEX for the flat ablation arm, not a fourth role. The slate band that used to sit
% here was assembly_introduced from the atlas origin palette, which is an atom origin recovered by
% atom-mapped forward replay after generation, not a coordinate P assigns before sampling. On this
% axis it asserted a fourth precursor role in a three-component reaction.
% Equal-weight bands. The role colours differ widely in intrinsic luminance -- the orange is far
% lighter than the others -- so one shared opacity made isocyanide read as the faintest stratum,
% which asserts an importance ordering the roles do not have. The per-band fill opacity, rule
% opacity and label mix were solved so every band sits the same L* below the ground and every label
% sits at the same L*. Hue separates the roles; value deliberately does not.
% Solved against groundcore, the vignette's centre value, which is what sits under the middle of
% the block. The vignette fades outward, so the equalisation is exact at the centre and drifts
% toward the block's two ends, where no atom is drawn. Deepening the vignette widens that drift:
% these numbers must be re-solved if groundcore changes again.
% The first column is the stratum offset in mm inside a state scope, which is exactly what
% \forgesites uses to place atoms. Bands and atoms therefore read from one set of numbers.
\foreach \sy/\col/\bandop/\ruleop/\labmix/\lb/\grp in {%
       8.4/roleamine/0.082/0.37/83/{amine}/{$\mathrm{N}\text{--}\mathrm{H}$},
       0.0/roleald/0.087/0.39/75/{aldehyde}/{$\mathrm{CH}{=}\mathrm{O}$},
      -8.4/roleiso/0.156/0.70/57/{isocyanide}/{$\mathrm{N}{\equiv}\mathrm{C}$}} {
  \pgfmathsetmacro\yc{\ya+\sy/10}
  \pgfmathsetmacro\ytop{\yc+\bandh}
  \pgfmathsetmacro\ybot{\yc-\bandh}
  \fill[\col, opacity=\bandop] (\bandl,\ybot) rectangle (\bandr,\ytop);
  \draw[\col, opacity=\ruleop, line width=0.4pt] (\bandl,\ybot) -- (\bandr,\ybot);
  \node[font=\tiny, text=\col!\labmix!black, anchor=west] at (\labx,\yc+0.10) {\lb};
  % the group each role contributes at the reaction core, quoted from the reaction SMARTS
  \node[font=\tiny, text=inkmuted!70, anchor=west] at (\labx,\yc-0.11) {\grp};
}
\pgfmathsetmacro\blocktop{\ya+0.84+\bandh+0.06}
\pgfmathsetmacro\blockbot{\ya-0.84-\bandh-0.06}
\pgfmathsetmacro\rolehdr{\blocktop+0.10}
\draw[draw=inkmuted!22, line width=0.35pt] (\bandl,\blocktop) -- (\bandr,\blocktop);
\draw[draw=inkmuted!22, line width=0.35pt] (\bandl,\blockbot) -- (\bandr,\blockbot);
\node[font=\tiny, text=inkmuted!75, anchor=west] at (\labx,\rolehdr) {role $r_i$};

%% ================================================================ endpoints
\fill[peri] (\inx,\ya) circle (1.1mm);
\node[font=\small, text=peri!85!black, anchor=south] at (\inx,\ya+0.16) {$P$};
\node[font=\tiny, text=inkmuted, anchor=north] at (\inx,\ya-0.16) {program in};
\fill[cyanacc] (\outx,\ya) circle (1.1mm);
\node[font=\small, text=cyanacc!80!black, anchor=south] at (\outx,\ya+0.16) {$T$};
\node[font=\tiny, text=inkmuted, anchor=north] at (\outx,\ya-0.16) {route out};

%% ================================================================ the three states
% One shared position layout across all three: the reaction-program coordinates are per-position and
% are given by P before sampling starts, so every position keeps its role colour and its core halo
% throughout. Opacity, not colour, carries what the flow actually determines.
\newcommand{\forgesites}[1]{%
  % x is position along the chain, y is the role stratum, so a site's height IS its role coordinate
  \foreach \px/\py/\col/\core in {%
      -7.4/8.4/roleamine/1,
      -5.0/0.0/roleald/1, -2.0/0.0/roleald/0, 1.0/0.0/roleald/1,
       4.0/0.0/roleald/0,  7.0/0.0/roleald/0,
       2.9/-8.4/roleiso/1,  5.4/-8.4/roleiso/0} {
    \ifnum\core=1 \draw[\col, line width=0.45pt, opacity=0.42] (\px mm,\py mm) circle (1.7mm); \fi
    \ifnum#1=1 \fill[\col] (\px mm,\py mm) circle (0.95mm);
    \else \fill[\col, opacity=0.38] (\px mm,\py mm) circle (0.95mm); \fi
  }%
}

\begin{scope}[shift={(4.05,\ya)}] \forgesites{0} \end{scope}

\begin{scope}[shift={(\cx,\ya)}]
  \draw[roleald, line width=1.15pt] (-5mm,0mm) -- (-2mm,0mm) -- (1mm,0mm);
  \forgesites{0}
  \foreach \px in {-5.0, -2.0, 1.0} \fill[roleald] (\px mm,0mm) circle (0.95mm);
\end{scope}

\begin{scope}[shift={(9.95,\ya)}]
  \draw[roleald, line width=1.15pt] (-5mm,0mm) -- (7mm,0mm);
  \draw[roleiso, line width=1.15pt] (1mm,0mm) -- (2.9mm,-8.4mm) -- (5.4mm,-8.4mm);
  \draw[roleamine, line width=1.15pt] (-5mm,0mm) -- (-7.4mm,8.4mm);
  \forgesites{1}
\end{scope}

%% ================================================================ sampling-time axis
\draw[draw=inkmuted!26, line width=0.4pt] (\bandl,\yaxis) -- (\bandr,\yaxis);
\foreach \tx/\tl/\sl in {4.05/{$t=0$}/{$X_0$}, \cx/{$t$}/{$X_t$}, 9.95/{$t=1$}/{$X_1\!\to\!G$}} {
  \draw[draw=inkmuted!34, line width=0.5pt] (\tx,\yaxis) -- (\tx,\yaxis-0.14);
  \node[font=\tiny, text=inkmuted!85, anchor=north] at (\tx,\yaxis-0.17) {\tl};
  \node[statelab, anchor=north] at (\tx,\yaxis-0.47) {\sl};
}

%% ================================================================ mask chips, one each side
\node[font=\tiny, text=inkmuted, anchor=south] (chL) at (4.05,2.80)
  {$r_i$, $k_i$ and the core states from $P$: never resampled};
\node[font=\tiny, text=inkmuted, anchor=south] (chR) at (9.95,2.80)
  {atom and bond states outside the core: resampled};

%% ================================================================ coordinate panel, centred
% Left half: one column per graph position, in the same order and colours as the sites above, with
% the outermost pair bracketed to their atoms. Right half: the same coordinates entering the
% per-position embedding. Together they are what "reaction-program coordinates" means.
%% ================================================================ the core the reaction forms
% Quoted from the product side of the atom-mapped Ugi-3CR SMARTS,
%   ... >> [N:1][CH1:2][C+0:3](=O)[NH1+0:4]
% so this is N-CH-C(=O)-NH: the five-atom core the adapter owns. Each atom carries the colour of
% the component it came from, which is why the carbonyl O is grey -- map_1..map_4 come from the
% three precursors, but the O is template_introduced_0 and enters from no precursor at all. Halos
% mark core membership, the same convention the states above use for k_i.
% Centred on \cx, under no single state, because the core belongs to all three. The adapter owns
% this chemistry and P fixes it: "the chemistry adapter marks any atom, pointer, or bond states
% fixed by P. FORGE holds these states constant while corrupting the remaining states." So the
% core is already there at t=0 and is never sampled. Hanging it under X_1 said the opposite --
% that it emerges at the end -- which is the reverse of what the figure argues.
% The two open bonds mean "continues into the precursor region", not "exactly one R": the amine
% is [NX3;H2,H1], so N:1 may carry one substituent or two, and the drawing does not commit.
\begin{scope}[shift={(6.49,-1.60)}]
  \draw[draw=inkmuted!55, line width=0.5pt]
    (-3.4mm,1.7mm) -- (0mm,0mm) -- (3.4mm,1.7mm) -- (6.8mm,0mm) -- (10.2mm,1.7mm) -- (13.6mm,0mm);
  \draw[draw=inkmuted!55, line width=0.5pt] (6.45mm,0mm) -- (6.45mm,-3.4mm);
  \draw[draw=inkmuted!55, line width=0.5pt] (7.15mm,0mm) -- (7.15mm,-3.4mm);
  \foreach \ax/\ay/\col in {0/0/roleamine, 3.4/1.7/roleald, 6.8/0/roleiso,
                            10.2/1.7/roleiso, 6.8/-3.4/inkmuted} {
    \draw[\col, line width=0.4pt, opacity=0.42] (\ax mm,\ay mm) circle (1.5mm);
    \fill[\col] (\ax mm,\ay mm) circle (0.85mm);
  }
  \node[font=\tiny, text=roleamine!83!black, anchor=north] at (0mm,-1.6mm) {N};
  \node[font=\tiny, text=roleiso!57!black, anchor=west] at (12.0mm,1.7mm) {N};
  \node[font=\tiny, text=inkmuted, anchor=north] at (6.8mm,-4.6mm) {O};
\end{scope}
% One tag only. Without it the inset is an unlabelled fragment, but the reason the carbonyl is
% grey is prose, and prose belongs in the figure caption where it costs no space in the drawing.
\node[font=\tiny, text=inkmuted!85, anchor=west] at (8.15,-1.60)
  {reaction core, marked by $k_i$};

\begin{scope}[shift={(2.88,\ypanel)}]
  \node[font=\tiny, text=ink, anchor=east] at (-0.6mm,3.05mm) {$r_i$};
  \node[font=\tiny, text=ink, anchor=east] at (-0.6mm,0.05mm) {$k_i$};
  \foreach \i/\col/\core in {0/roleamine/1, 1/roleald/1, 2/roleald/0, 3/roleald/1,
                             4/roleald/0, 5/roleald/0, 6/roleiso/1, 7/roleiso/0} {
    \pgfmathsetmacro\x{\i*3.0}
    \fill[\col] (\x mm,2.2mm) rectangle (\x mm+2.4mm,4.6mm);
    \ifnum\core=1
      \fill[peri] (\x mm,-0.8mm) rectangle (\x mm+2.4mm,1.6mm);
    \else
      \draw[peri!55, line width=0.4pt] (\x mm,-0.8mm) rectangle (\x mm+2.4mm,1.6mm);
    \fi
  }
\end{scope}
\foreach \ax/\i in {-7.4/0, 7.0/7} {
  \pgfmathsetmacro\axcm{4.05+\ax/10}
  \pgfmathsetmacro\colx{2.88+\i*0.30+0.12}
  \draw[draw=inkmuted!30, line width=0.35pt] (\axcm,-1.26) -- (\colx,\ypanel+0.50);
}

\begin{scope}[shift={(8.05,\ypanel)}]
  \foreach \i/\lb/\op in {0/{$e_c(c)$}/0.85, 1/{$e_d(d)$}/0.68, 2/{$e_r(r_i)$}/0.50, 3/{$e_k(k_i)$}/0.34} {
    \pgfmathsetmacro\x{\i*7.8}
    \fill[peri, opacity=\op] (\x mm,1.4mm) rectangle (\x mm+5.2mm,4.0mm);
    \node[font=\tiny, text=inkmuted, anchor=north] at (\x mm+2.6mm,0.9mm) {\lb};
    \ifnum\i<3 \node[font=\tiny, text=ink] at (\x mm+6.5mm,2.7mm) {$+$}; \fi
  }
  \draw[-{Stealth[length=3.5pt,width=3.5pt]}, draw=peri, line width=0.8pt] (28.8mm,2.7mm) -- (32.2mm,2.7mm);
  \node[font=\tiny, text=inkmuted, anchor=south] at (30.5mm,3.9mm) {LN};
  \fill[navybar] (33.0mm,1.4mm) rectangle (39.4mm,4.0mm);
  \node[font=\tiny, text=inkmuted, anchor=north] at (36.2mm,0.9mm) {$h_i^{P}$};
\end{scope}

%% ================================================================ equation plates, centred
\node[plate, anchor=south] at (\cx,3.42)
  {{\plateref sampling step}\\[2pt]
   $X_t=(\mathbf{a},\boldsymbol{\pi},\mathbf{b},\mathbf{e}^{L},\mathbf{e}^{R},\mathbf{z})
    \qquad
    X_{k+1}=\Pi_P\!\left(\Phi_k^\theta(X_k,\xi_k)\right)$\\[3pt]
   {\plateref $\Pi_P$ restores every coordinate $P$ fixes; the rest keep their sampled value}};
\node[plate, anchor=north] at (\cx,-3.62)
  {{\plateref acceptance and route}\\[2pt]
   $\mathrm{L1}(G,P)=\mathbf{1}\!\left[\widehat{\mathcal{E}}_P(G)\neq\varnothing\right]
    \qquad T=\mathcal{A}(G,P)\in\mathfrak{T}\cup\{\bot\}$};

\end{tikzpicture}
